\documentclass{article}
\usepackage[utf8]{inputenc}
\usepackage[T1]{fontenc}
\usepackage{geometry}
\geometry{
    a4paper,
    left=2.5cm,
    right=2.5cm,
    top=2.5cm,
    bottom=2.5cm
}
\usepackage{graphicx}
\usepackage{float}
\usepackage{hyperref}
\usepackage[english]{babel}
\usepackage{amsmath}
\usepackage{amssymb}
\usepackage{array}
\usepackage{booktabs}
\usepackage{caption}
\usepackage{subcaption}
\usepackage{enumitem}
\usepackage{fancyhdr}
\usepackage{titlesec}
\usepackage{url}
\usepackage{textcomp}
\usepackage{xcolor}
\usepackage{listings}
\usepackage{multirow}
\usepackage{longtable}
\usepackage{tabularx}
\usepackage{times}

% Set up page style
\pagestyle{fancy}
\fancyhf{}
\rfoot{\thepage}
\renewcommand{\headrulewidth}{0pt}

% Title formatting
\titleformat{\chapter}[display]{\normalfont\bfseries}{\chaptertitle\thechapter}{0pt}{\Huge}
\titlespacing*{\chapter}{0pt}{50pt}{40pt}

% Document info
\title{Template Document}
\author{Academic Support Agent}
\date{\today}

\begin{document}
\maketitle

%%%%%%%%%%%%%%%%%%%%%%%
		%%  Slide 1: TITOLO  %%
		%%%%%%%%%%%%%%%%%%%%%%%
\begin{frame}
\setTitlestyleDissertation
\maketitle
\end{frame}

		%%%%%%%%%%%%%%%%%%%%%%%%%%%%%%%%%%%%%
		%%  Slide 2: <ArgomentoPrimaSlide> %%
		%%%%%%%%%%%%%%%%%%%%%%%%%%%%%%%%%%%%%

\section{Capitolo 1}
\begin{frame}
    \frametitle{Titolo prima slide}
    Elenco puntato: 
    \begin{itemize}
        \item Primo elemento
        \item<alert@1> Secondo elemento evidenziato 
        \item Terzo elemento
            \begin{itemize}
                \item Secondo layer dell'elenco puntato
                \item ... 
            \end{itemize}
    \end{itemize}
\end{frame}

		%%%%%%%%%%%%%%%%%%%%%%%%%%%%%%%%%%%%%%%%
		%%  Slide 3: <ArgomentoSecondaSlide>  %%
		%%%%%%%%%%%%%%%%%%%%%%%%%%%%%%%%%%%%%%%%

\section{Capitolo 2}
\begin{frame}
    \frametitle{Titolo seconda slide}
    Elenchi numerati 
    \begin{enumerate}
        \item primo elemento 
        \item secondo elemento
            \begin{enumerate}
                \item Esempio di struttura annidata
                \item ...
            \end{enumerate}
    \end{enumerate}
\end{frame}



        %%%%%%%%%%%%%%%%%%%%%%%%%%%%%%%%%%%%
        %% Slide 4: <ArgomentoTerzaSlide> %%
        %%%%%%%%%%%%%%%%%%%%%%%%%%%%%%%%%%%%

\section{Capitolo 3}
\begin{frame}
    \frametitle{Titolo terza slide}
    contenuto terza slide
\end{frame}
\end{document}